% =============================================================================
% HISTORICAL AND FORENSIC GLOSSARY — SUPPORT APPENDIX FOR THE READER
% =============================================================================
\cleardoublepage
\thispagestyle{empty}
\chapter*{Historical and Forensic Glossary}
\addcontentsline{toc}{chapter}{Historical and Forensic Glossary}
\markboth{Historical and Forensic Glossary}{}

\begingroup

\noindent\textit{Support vocabulary for navigating the documents of the 1915 drought and the psychiatric reports of the Barbacena Hospital-Colony (1917--1979).}

\vspace{0.8em}
\noindent\rule{\textwidth}{0.8pt}
\vspace{0.6em}

\noindent\textbf{1. TERMS OF THE DROUGHT AND THE SERTÃO (1914--1915)}

\vspace{0.4em}
\begin{itemize}[itemsep=0.2em]
  \item \textbf{Alagadiço Camp (1915):} Camp for the confinement of retirantes created near Fortaleza during the 1915 drought. It is the historical referent adopted for the fictional setting in the first part of this book.
  \item \textbf{Patu Camp (1932--1933):} Camp located in Senador Pompeu and associated with the 1932 drought. It is not the 1915 camp in Joaquim's route.
  \item \textbf{Retirantes / Molambudos:} Popular, pejorative designation for the sertanejos who fled the interior on foot, in rags of cloth (molambos), trying to survive the extreme drought.
  \item \textbf{Mass Grave:} Shallow collective pit where hundreds of bodies of famine and typhus victims were deposited daily, unidentified, in the concentration camps of 1915.
\end{itemize}

\vspace{0.8em}
\noindent\rule{\textwidth}{0.5pt}
\vspace{0.6em}

\noindent\textbf{2. MEDICAL AND PSYCHIATRIC TERMS OF THE TWENTIETH CENTURY}

\vspace{0.4em}
\begin{itemize}[itemsep=0.2em]
  \item \textbf{Barbacena Hospital-Colony:} Psychiatric institution founded in Minas Gerais in 1903. Between 1930 and 1980, it became a center of mass internment of marginalized people without formal psychiatric diagnosis, resulting in the death of more than 60,000 people (an episode known as the \textit{Brazilian Holocaust}).
  \item \textbf{Electroconvulsive Therapy (ECT / Electroshock):} Application of electric current to the brain to induce seizures. At the Hospital-Colony, it was practiced continuously and punitively, without anesthesia or conductive gel, causing burns and fractures of the femur and spine.
  \item \textbf{Transorbital Lobotomy (Freeman Procedure):} Surgical intervention introduced in the 1940s in which a metal rod (similar to an ice pick) was inserted through the inner corner of the eye into the brain to sever connections of the frontal lobe.
  \item \textbf{Cachexia:} State of extreme chronic malnutrition and profound muscular wasting.
  \item \textbf{Beau's Lines:} Transverse grooves in the nails caused by the temporary interruption of nail growth during severe episodes of hunger or infection.
  \item \textbf{Unilateral Jaundice:} Yellowish coloration restricted to a single eye, an anomalous forensic finding documented in the examination of Joaquim (Patient 1,260).
\end{itemize}

\vspace{0.8em}
\noindent\rule{\textwidth}{0.5pt}
\vspace{0.6em}

\noindent\textbf{3. FORENSIC AND ARCHIVAL TERMS (2026)}

\vspace{0.4em}
\begin{itemize}[itemsep=0.2em]
  \item \textbf{Leucotome:} Pointed metal surgical instrument used in lobotomies.
  \item \textbf{Cribriform Plate:} Perforated bone plate at the roof of the nasal/ocular cavity that provides access to the base of the brain.
  \item \textbf{Ergodic Fictional Pact:} Narrative device in which the physical structure of the work and the reader's decisions simulate the very contagion described in the diary.
\end{itemize}

\endgroup
\clearpage
